\chapter{Ground spaces, gaps and Hermitian terms}
\label{ch:hamiltonians}

This chapter treats the spectral notions attached to the Hamiltonian of an
interaction, with the notation of Chapters~\ref{ch:sites}
and~\ref{ch:dynamics}. The ground space and the gap are defined for any
square matrix, so they apply to the Hamiltonian of an interaction as it stands.
The gap hypothesis of the two-dimensional area law~\cite{OpenAI2026AreaLaw} and
of the polynomial approximation by projected entangled pair
states~\cite{OpenAI2026PolynomialPEPS} is shown to be the same as a gap with a
one-dimensional ground space.

\section{Ground spaces and gaps}
\label{sec:hamiltonians_gap}

\begin{definition}[Ground space and gap]
    \label{def:hamiltonians_gap}
    \lean{ManyBody.groundSpace, ManyBody.mem_groundSpace_iff,
        ManyBody.mem_groundSpace_iff_mulVec, ManyBody.IsGapped}
    \leanok
    Let $H$ be a complex $n\times n$ matrix and $E_0\in\R$. The ground space
    of $H$ at $E_0$ is the eigenspace $\ker(H-E_0)$ of $\C^n$. For $\Delta\in\R$,
    the matrix $H$ is gapped at $E_0$ with gap $\Delta$ when $\ker(H-E_0)\neq0$
    and every eigenvalue $\mu\neq E_0$ of $H$ is real with $\mu\ge E_0+\Delta$.
    For $\Delta>0$, this makes $E_0$ the least eigenvalue and $\Delta$ a lower
    bound for the distance from $E_0$ to the rest of the spectrum.
\end{definition}

\begin{lemma}[The gap in terms of the eigenvalues]
    \label{lem:hamiltonians_gap_eigenvalues}
    \lean{Matrix.IsHermitian.hasEigenvalue_toEuclideanLin_iff,
        ManyBody.isGapped_iff_eigenvalues}
    \leanok
    \uses{def:hamiltonians_gap}
    Let $H$ be Hermitian with eigenvalues $\lambda_1,\dots,\lambda_n$, counted
    with multiplicity. Then $H$ is gapped at $E_0$ with gap $\Delta$ if and only
    if $\ker(H-E_0)\neq0$ and every $\lambda_k$ equals $E_0$ or is at least
    $E_0+\Delta$.
\end{lemma}
\begin{proof}\leanok
    The eigenvalues of $H$ are its points of spectrum, which are the
    $\lambda_k$.
\end{proof}

\begin{lemma}[The operator gap inequality]
    \label{lem:hamiltonians_gap_operator_inequality}
    \lean{ManyBody.IsGapped.posSemidef_sub_smul_one,
        ManyBody.mulVec_eq_of_groundSpace_eq_span}
    \leanok
    \uses{def:hamiltonians_gap, lem:hamiltonians_gap_eigenvalues}
    Let $H$ be Hermitian, let $\Omega\in\C^n$ with $\lVert\Omega\rVert=1$, and let
    $\Delta\in\R$. If $H$ is gapped at $E_0$ with gap $\Delta$ and
    $\ker(H-E_0)=\C\Omega$, then $H\Omega=E_0\Omega$ and
    $H-E_0\ge\Delta(1-\lvert\Omega\rangle\langle\Omega\rvert)$. This is the gap
    hypothesis of the two-dimensional area law~\cite{OpenAI2026AreaLaw}
    (\texttt{00-introduction.tex}, lines 44--47).
\end{lemma}
\begin{proof}\leanok
    Write $G=H-E_0-\Delta(1-\lvert\Omega\rangle\langle\Omega\rvert)$, and let
    $b_1,\dots,b_n$ be an orthonormal basis of eigenvectors of $H$, with
    $Hb_k=\lambda_kb_k$. If $\lambda_k=E_0$, then $b_k\in\C\Omega$, so
    $b_k=\langle\Omega,b_k\rangle\Omega$ and $Gb_k=0$. Otherwise
    $\lambda_k\ge E_0+\Delta$ by Lemma~\ref{lem:hamiltonians_gap_eigenvalues}, and
    $b_k\perp\Omega$, so $Gb_k=(\lambda_k-E_0-\Delta)b_k$. Thus $G$ is Hermitian with an
    orthonormal basis of eigenvectors with nonnegative eigenvalues, and
    $\langle\psi,G\psi\rangle=\sum_k\nu_k\lvert\langle b_k,\psi\rangle\rvert^2\ge0$ for
    every $\psi$.
\end{proof}

\begin{theorem}[A gapped unique ground state]
    \label{thm:hamiltonians_gapped_unique}
    \lean{ManyBody.isGapped_and_groundSpace_eq_span_iff}
    \leanok
    \uses{def:hamiltonians_gap}
    Let $\Delta>0$ and let $\Omega\in\C^n$. The following are equivalent.
    \begin{enumerate}
        \item $\lVert\Omega\rVert=1$, $H\Omega=E_0\Omega$ and
              $H-E_0\ge\Delta(1-\lvert\Omega\rangle\langle\Omega\rvert)$.
        \item $H$ is Hermitian, $\lVert\Omega\rVert=1$, $H$ is gapped at $E_0$
              with gap $\Delta$, and $\ker(H-E_0)=\C\Omega$.
    \end{enumerate}
\end{theorem}
\begin{proof}\leanok
    \uses{thm:positive_gap_uniqueness, lem:hamiltonians_gap_operator_inequality}
    Write $G=H-E_0-\Delta(1-\lvert\Omega\rangle\langle\Omega\rvert)$.

    Assume (1). The matrix $H$ is $G$ plus Hermitian matrices, hence Hermitian,
    and $\ker(H-E_0)=\C\Omega$ by
    Theorem~\ref{thm:positive_gap_uniqueness}. Let $\psi\neq0$ satisfy
    $H\psi=\mu\psi$ with $\mu\neq E_0$. Then $\mu$ is real, and
    $\langle\Omega,\psi\rangle=0$ because $\Omega$ is an eigenvector of the
    Hermitian matrix $H$ at a different eigenvalue. Hence
    $0\le\langle\psi,G\psi\rangle=(\mu-E_0-\Delta)\lVert\psi\rVert^2$, so
    $\mu\ge E_0+\Delta$.

    Assume (2). Then (1) is
    Lemma~\ref{lem:hamiltonians_gap_operator_inequality}.
\end{proof}

\section{Hermitian parts of the terms}
\label{sec:hamiltonians_hermitian_part}

\begin{definition}[Hermitian parts of the terms]
    \label{def:hamiltonians_hermitian_part}
    \lean{ManyBody.Interaction.hermitianPart, ManyBody.Interaction.hermitianPart_term}
    \leanok
    \uses{def:dynamics_interaction, thm:sites_supported_subalgebra}
    For an interaction $(h_i)_{i\in I}$ with supports $(S_i)_{i\in I}$, the
    interaction of its Hermitian parts has the terms
    $\operatorname{Re}h_i=\tfrac12(h_i+h_i^*)$ and the same supports. The term
    $\operatorname{Re}h_i$ is supported on $S_i$, since the operators supported
    on a region are closed under adjoints and real linear
    combinations~\cite{OpenAI2026PolynomialPEPS}.
\end{definition}

\begin{lemma}[Hermitian parts preserve a Hermitian Hamiltonian]
    \label{lem:hamiltonians_hermitian_part}
    \lean{ManyBody.Interaction.isHermitian_hermitianPart_term,
        ManyBody.Interaction.norm_hermitianPart_term_le,
        ManyBody.Interaction.hamiltonian_hermitianPart}
    \leanok
    \uses{def:hamiltonians_hermitian_part}
    Each $\operatorname{Re}h_i$ is Hermitian with
    $\lVert\operatorname{Re}h_i\rVert\le\lVert h_i\rVert$. If the index set is
    finite and $\sum_ih_i$ is Hermitian, then
    $\sum_i\operatorname{Re}h_i=\sum_ih_i$.
\end{lemma}
\begin{proof}\leanok
    The norm bound is the triangle inequality for
    $\tfrac12(h_i+h_i^*)$. The real part is real linear, so
    $\sum_i\operatorname{Re}h_i=\operatorname{Re}\sum_ih_i$, which is
    $\sum_ih_i$ when the sum is Hermitian.
\end{proof}

\section{Lattice models}
\label{sec:hamiltonians_models}

\begin{definition}[Lattice models]
    \label{def:hamiltonians_models}
    \lean{ManyBody.ChainInteraction, ManyBody.DomainInteraction,
        ManyBody.PlaneInteraction, ManyBody.RectangleInteraction,
        ManyBody.SquareInteraction}
    \leanok
    \uses{def:dynamics_interaction, def:sites_chain_support, def:sites_line_domain,
        def:sites_plane_domain, def:sites_square_lattice}
    Each lattice model is an interaction with local dimension $q$ at every site,
    for a choice of sites and supports: the chain of $L$ sites with its supports
    of range $R$; a finite set $\Lambda\subset\Z$ with its supports of range $R$;
    a finite set $\Lambda\subset\Z^2$ with its supports of range
    $R$~\cite{OpenAI2026AreaLaw}; and the $w\times h$ square lattice with one term
    on each vertex and one on the two endpoints of each
    edge~\cite{OpenAI2026PolynomialPEPS}, of which the $L\times L$ square is the
    case $w=h=L$. Hermiticity and norm bounds of the terms are hypotheses of the
    statements about a model, not part of the model.
\end{definition}
